% !TEX TS-program = pdflatex
% CV ciblé — Fullstack Developer (TypeScript) — Shindan
\documentclass[a4paper,10pt]{article}
\usepackage[utf8]{inputenc}
\usepackage[T1]{fontenc}
\usepackage[scaled=0.94]{helvet}
\renewcommand{\familydefault}{\sfdefault}
\usepackage[left=1.2cm,right=1.2cm,top=0.7cm,bottom=0.6cm]{geometry}
\usepackage{enumitem}
\usepackage{titlesec}
\usepackage{xcolor}
\usepackage[hidelinks]{hyperref}
\definecolor{navy}{RGB}{23,54,93}
\pagestyle{empty}
\setlength{\parindent}{0pt}
\setlength{\parskip}{0pt}
\linespread{0.86}
\hypersetup{pdftitle={CV - Baha Eddine Belhaj Mouldi - Fullstack Developer TypeScript},pdfauthor={Baha Eddine Belhaj Mouldi},pdfsubject={Fullstack, TypeScript, APIs, Sécurité Mobile}}
\titleformat{\section}{\normalsize\bfseries\color{navy}}{}{0pt}{\MakeUppercase}[{\vspace{1pt}\color{navy}\hrule height 0.6pt}]
\titlespacing*{\section}{0pt}{3pt}{1.5pt}
\setlist[itemize]{leftmargin=1.1em,label=\textbullet,itemsep=0.1pt,topsep=0.2pt,parsep=0pt}
\newcommand{\cventry}[4]{\textbf{#1} \hfill \textbf{#4}\\[1pt]#2{\ifx\relax#3\relax\else, #3\fi}\par\vspace{1pt}}
\begin{document}
\begin{center}
{\LARGE\bfseries\color{navy} Baha Eddine Belhaj Mouldi}\\[3pt]
{\large Fullstack Developer \textbar\ TypeScript \textbar\ Sécurité Mobile}\\[5pt]
Tunis, Tunisie \quad|\quad +216 55 128 982 \quad|\quad \href{mailto:bahaeddine.belhajmouldi@gmail.com}{bahaeddine.belhajmouldi@gmail.com}\\[1pt]
\href{https://www.linkedin.com/in/bahamouldi}{linkedin.com/in/bahamouldi} \quad|\quad \href{https://github.com/bahamouldi}{github.com/bahamouldi} \quad|\quad \href{https://bahamouldi.github.io/portfolio/}{bahamouldi.github.io/portfolio}
\end{center}
\vspace{1pt}
\section{Profil}
Ingénieur en génie informatique orienté développement fullstack et sécurité applicative. Expérience concrète en développement backend/frontend TypeScript, conception d'APIs REST, développement mobile cross-platform (Flutter/Android/iOS) et analyse de sécurité d'applications mobiles (reverse engineering). Passionné par la qualité logicielle : tests unitaires, revues de code et architecture maintenable. Recherche une alternance Fullstack Developer pour contribuer à des produits innovants alliant web, mobile et sécurité.
\section{Compétences clés}
\textbf{Backend \& APIs} : TypeScript, Node.js, Next.js, FastAPI, Python, Java, Spring Boot, REST APIs, architecture microservices\\[0.5pt]
\textbf{Frontend \& Mobile} : React, Next.js, Tailwind CSS, Flutter, Dart, Android/iOS, interfaces responsive mobile-first\\[0.5pt]
\textbf{Données \& ORM} : PostgreSQL, MySQL, MongoDB, Prisma, SQL, NoSQL\\[0.5pt]
\textbf{DevOps \& Outils} : Docker, Kubernetes, Git, GitHub Actions, CI/CD, Linux, Bash\\[0.5pt]
\textbf{Sécurité} : Reverse engineering d'applications mobiles, OWASP Top 10, tests d'intrusion web/API, analyse de vulnérabilités, codage sécurisé\\[0.5pt]
\textbf{Qualité} : Tests unitaires, tests d'intégration, revue de code, architecture modulaire, documentation technique
\section{Expérience professionnelle}
\cventry{Ingénieur Sécurité Logiciel --- Stagiaire PFE}{Beehive Entreprises}{Tunisie}{01/2026 -- 05/2026}
\begin{itemize}
  \item Développement backend (FastAPI, Python) de BeeWAF : modules anti-bot, DLP, anti-DDoS, virtual patching (37 CVE).
  \item Architecture Docker/Kubernetes, reporting ELK, détection ML d'attaques web. Score 98,2/100, 0\% faux positifs.
  \item Rédaction de documentation technique et intégration de logique de sécurité dans un prototype déployable.
\end{itemize}
\cventry{Spécialiste Cybersécurité Junior, Freelance}{Beehive Entreprises}{Tunisie}{01/2025 -- 05/2026}
\begin{itemize}
  \item Tests d'intrusion sur 8 applications web et API, production de rapports de remédiation détaillés.
  \item Reverse engineering d'applications mobiles (Android/iOS) pour analyse de sécurité et détection de failles.
\end{itemize}
\cventry{Chef de projet technique — Vmate}{ENIT}{Tunisie}{06/2024}
\begin{itemize}
  \item Encadrement d'une équipe de 4 sur une plateforme sécurisée web/mobile (projet classé 1er).
  \item Développement sécurisé, revue de code et tests de sécurité applicative.
\end{itemize}
\cventry{Stagiaire Cybersécurité — Détection menaces SAP}{Streamlink}{Tunisie}{07/2025 -- 08/2025}
\begin{itemize}
  \item Automatisation de la collecte d'événements via Python/PyRFC et support de corrélation de sécurité ELK.
\end{itemize}
\section{Projets Fullstack}
\cventry{Fedi Phone — Plateforme E-commerce}{Next.js, TypeScript, Prisma, PostgreSQL, Tailwind CSS, Docker}{}{2026}
\begin{itemize}
  \item Application e-commerce full-stack : authentification JWT, panier, wishlist, recherche temps réel, dashboard admin (CRUD produits/commandes/utilisateurs), intégration WhatsApp.
  \item Architecture frontend/backend TypeScript, ORM Prisma, design responsive mobile-first, mode sombre, animations Framer Motion.
\end{itemize}
\cventry{Architecture Microservices}{TypeScript, Node.js, Microservices, REST API}{}{2026}
\begin{itemize}
  \item Application distribuée en microservices : services découplés, communication REST, organisation modulaire orientée scalabilité et maintenabilité.
\end{itemize}
\cventry{Oumaya — Application Mobile}{Flutter, Dart, Android, iOS}{}{2025}
\begin{itemize}
  \item Application mobile cross-platform (Android/iOS) avec interface soignée et logique applicative structurée.
\end{itemize}
\cventry{Système de Gestion de Cinéma}{Java, Spring Boot, JPA, MySQL, REST API}{}{2025}
\begin{itemize}
  \item API REST backend : gestion de films, séances, réservations, paiements et statistiques.
\end{itemize}
\cventry{BeeWAF — Pare-feu Applicatif Intelligent}{FastAPI, Python, Docker, Kubernetes, ELK}{}{2026}
\begin{itemize}
  \item Modules backend sécurité : anti-bot, DLP, anti-DDoS, virtual patching, reporting orienté monitoring.
\end{itemize}
\section{Formation}
\cventry{Diplôme d'Ingénieur en Génie Informatique}{ENIT}{Tunisie}{09/2023 -- 06/2026}
\cventry{Classes préparatoires (Physique-Technologie)}{IPEI El Manar}{Tunisie}{09/2021 -- 06/2023}
\section{Certifications}
Git \& GitHub --- 365 Data Science (2024) ; Réseaux --- Cisco (2025) ; Fondements du Deep Learning --- NVIDIA (2025) ; Machine Learning --- NVIDIA (2024) ; IA Adversariale --- NVIDIA (2025).
\section{Langues}
Anglais (B2) \quad|\quad Français (B2) \quad|\quad Arabe (Natif) \quad|\quad Chinois (Notions)
\end{document}
